% =============================================================================
% Module 14: Solutions to Exercises
% =============================================================================
% This file is conditionally included from main.tex via \ifsolutions
% Every number is attested by Solutions/14_Multiplane_DSPL/problems_14.wl
% \annot{check ids quoted in each [Attested: ...] line}.
% =============================================================================

\section{Solutions to Exercises}
\label{sec:mp_solutions}

\noindent
Full worked solutions are also available in the Mathematica script
\solnlink{14_Multiplane_DSPL/problems\_14.wl}{problems\_14.wl},
which verifies each answer symbolically and numerically.  Distances are
computed in the script from the comoving-distance integral of
eq.~\eqref{eq:mp_DA} for flat cosmologies, with
$\Hubble=70$\,km\,s$^{-1}$\,Mpc$^{-1}$ wherever an absolute distance is
quoted.  We write $\chi_x\equiv\chi(z_x)$.

% ----- Exercise 14.1 -----
\subsection*{Exercise~\ref{ex:mp_hubble_free}: $\Brat$ Is Hubble-Free}

The two observer-to-source distances start at $\chi(0)=0$; the two
lens-to-source distances start at $\chi_l$:
\begin{align*}
  D_{s1} &= \frac{c}{\Hubble}\,\frac{\tilde\chi_{s1}}{1+z_{s1}}, &
  D_{s2} &= \frac{c}{\Hubble}\,\frac{\tilde\chi_{s2}}{1+z_{s2}}, \\
  D_{l,s1} &= \frac{c}{\Hubble}\,\frac{\tilde\chi_{s1}-\tilde\chi_l}{1+z_{s1}}, &
  D_{l,s2} &= \frac{c}{\Hubble}\,\frac{\tilde\chi_{s2}-\tilde\chi_l}{1+z_{s2}},
\end{align*}
with $\tilde\chi=\int_0^z\dd z'/E$.  The $(1+z)$ in each denominator is
set by the far endpoint, so $D_{s1}$ and $D_{l,s1}$ both carry
$(1+z_{s1})^{-1}$, and $D_{s2}$ and $D_{l,s2}$ both carry
$(1+z_{s2})^{-1}$.  In
$\Brat=D_{l,s1}D_{s2}/(D_{s1}D_{l,s2})$ the numerator and the
denominator each contain two factors of $c/\Hubble$, one of
$(1+z_{s1})^{-1}$ and one of $(1+z_{s2})^{-1}$.  They cancel, leaving
\begin{equation*}
  \Brat = \frac{(\chi_{s1}-\chi_l)\,\chi_{s2}}{\chi_{s1}\,(\chi_{s2}-\chi_l)} .
\end{equation*}
The lens redshift never appears in a prefactor because no distance in
$\Brat$ ends at the lens; $z_l$ enters only through $\chi_l$.
Since $\Hubble$ multiplies every $\chi$ by the same constant, only the
shape of $E(z)$ survives.  Numerically, at the illustrative
geometry, $\Brat=0.6742$ for $\Hubble=50$, 70 and 100 alike.

\emph{[Attested: \solnlink{14_Multiplane_DSPL/problems\_14.wl}{problems\_14.wl},
checks 1a--1e.]}

% ----- Exercise 14.2 -----
\subsection*{Exercise~\ref{ex:mp_isothermal}: The Isothermal Ruler}

With $\gp=2$ the exponent in $\Brat=\etamea^{\gp-1}$ is one, so
$\Brat=\etamea$.  Directly: for $\gp=2$, eq.~\eqref{eq:mp_plaw} gives
$\alpha_i(\theta)=\thetaE^{(i)}$, a constant.  The pointwise relation
$\alpha_2=\alpha_1/\Brat$ then reads $\thetaE^{(2)}=\thetaE^{(1)}/\Brat$,
i.e.\ $\etamea=\Brat$.

The geometric reason is that an SIS bends every ray by the same physical angle
$\alphahat=4\pi\sigma_v^2/c^2$.  The reduced deflection on plane $i$ is
that angle times the efficiency, $\alpha_i=(D_{l,si}/D_{si})\alphahat$,
and for a constant-deflection lens the Einstein radius equals the
reduced deflection: the ray that lands on the ring is the one carried
exactly to the optical axis.  The deflector physics $\alphahat$ is common
to both planes and cancels:
\begin{equation*}
  \etamea=\frac{\alpha_1}{\alpha_2}
        =\frac{D_{l,s1}/D_{s1}}{D_{l,s2}/D_{s2}}=\Brat .
\end{equation*}
At $\gp=2$ the ratio of Einstein radii is the ratio of two lensing
efficiencies, with no residual dependence on $\sigma_v$.

\emph{[Attested: \solnlink{14_Multiplane_DSPL/problems\_14.wl}{problems\_14.wl},
checks 2a--2c.]}

% ----- Exercise 14.3 -----
\subsection*{Exercise~\ref{ex:mp_slope_lever}: The Slope Lever Arm}

\textbf{(a)} $\ln\Brat=(\gp-1)\ln\etamea$, so
\begin{equation*}
  \pderiv{\ln\Brat}{\gp}\bigg|_{\etamea}=\ln\etamea ,
\end{equation*}
which contains no $\gp$.  Since $\etamea<1$ it is negative: a steeper
fitted slope lowers $\Brat$.

\textbf{(b)} $\etamea=\Brat^{1/(\gp-1)}=0.6742^{1/0.9}=0.6453$ and
$\ln\etamea=-0.4381$.  A $1\%$ error requires
$|\delta\gp|=0.01/|\ln\etamea|=0.0228$.

\textbf{(c)} $\delta\Brat/\Brat=\ln\etamea\,\delta\gp=-0.4381\times0.03
=-1.31\%$.  With $S_w=-0.0330$, eq.~\eqref{eq:mp_propagate} gives, at
fixed $\Omegam$, $\delta w=(-0.0131)/(-0.0330)=+0.40$.  A slope bias of a few hundredths,
with a perfectly measured $\etamea$, is enough to move the inferred $w$
by several tenths at this geometry.

\emph{[Attested: \solnlink{14_Multiplane_DSPL/problems\_14.wl}{problems\_14.wl},
checks 3a--3f.]}

% ----- Exercise 14.4 -----
\subsection*{Exercise~\ref{ex:mp_naive_distance}: The Cost of $D_{ds}=D_s-D_d$}

\textbf{(a)} In a flat universe
$D_d=\chi_d/(1+z_d)$, $D_s=\chi_s/(1+z_s)$ and
$D_{ds}=(\chi_s-\chi_d)/(1+z_s)$.  Comoving distances subtract, but the
naive difference divides the two comoving pieces by different factors:
\begin{equation*}
  (D_s-D_d)-D_{ds}=\frac{\chi_d}{1+z_s}-\frac{\chi_d}{1+z_d}
  =\chi_d\Big[\frac{1}{1+z_s}-\frac{1}{1+z_d}\Big],
\end{equation*}
which is strictly negative for every $z_s>z_d>0$ (it vanishes only in the
trivial limits $z_d\to0$ or $z_d\to z_s$).

\textbf{(b)} $D_d=1259.08$\,Mpc, $D_s=1726.62$\,Mpc, $D_{ds}=1097.08$\,Mpc,
while $D_s-D_d=467.54$\,Mpc.  The naive value is wrong by $629.54$\,Mpc,
more than half of the true $D_{ds}$.

\textbf{(c)} At $(0.5,1.0,2.0)$ the naive ``ratio'' is $0.8782$ against
the true $\Brat=0.6742$, an error of $+30.3\%$.  Read through
eq.~\eqref{eq:mp_propagate} this would be a shift of order $8$ in $w$,
far outside the linear regime, and far larger than any physical signal
in $\Brat$.

\emph{[Attested: \solnlink{14_Multiplane_DSPL/problems\_14.wl}{problems\_14.wl},
checks 4a--4f.]}

% ----- Exercise 14.5 -----
\subsection*{Exercise~\ref{ex:mp_limits}: Limits of $\Brat$}

\textbf{(a)} Setting $\chi_{s1}=\chi_l$ in eq.~\eqref{eq:mp_Bchi} gives
$\Brat=0$: a source on the lens plane is not lensed at all, so the ratio
of efficiencies vanishes.  Setting $\chi_{s2}=\chi_{s1}$ gives $\Brat=1$:
two coincident sources have equal efficiencies and carry no
cosmological information.

\textbf{(b)}
$\partial\Brat/\partial\chi_{s2}=-\Brat\,\chi_l/[\chi_{s2}(\chi_{s2}-\chi_l)]$,
negative for $0<\chi_l<\chi_{s2}$.

\textbf{(c)} For flat $\Lambda$CDM with $\Omegam=0.3$,
$\chi(\infty)=3.305\,c/\Hubble$; the integrand $1/E$ falls as
$(1+z)^{-3/2}$, so the integral converges.  Hence
$\Brat(0.5,1.0,z_{s2}\to\infty)=0.4943$: even an infinitely distant
second source gives a finite ratio, below the $0.6742$ of
$z_{s2}=2$.

\emph{[Attested: \solnlink{14_Multiplane_DSPL/problems\_14.wl}{problems\_14.wl},
checks 5a--5d.]}

% ----- Exercise 14.6 -----
\subsection*{Exercise~\ref{ex:mp_coupling}: How Big Is Lens--Lens Coupling?}

\textbf{(a)} An SIS of dispersion $\sigma$ deflects every ray by
$4\pi\sigma^2/c^2$.  Reduced to the far-source plane, the main lens gives
$\theta_{\rm E}(l\to s2)=4\pi(\sigma_L^2/c^2)\,D_{l,s2}/D_{s2}$ and the
near source gives
$\theta_{\rm E}(s1\to s2)=4\pi(\sigma_1^2/c^2)\,D_{s1,s2}/D_{s2}$.
Their ratio is
$(\sigma_1/\sigma_L)^2\,D_{s1,s2}/D_{l,s2}$; in a flat universe both
distances end at $z_{s2}$, so the $(1+z_{s2})$ factors cancel and the
ratio is $(\sigma_1/\sigma_L)^2(\chi_{s2}-\chi_{s1})/(\chi_{s2}-\chi_l)$.

\textbf{(b)} As $\chi_{s1}\to\chi_{s2}$ the numerator vanishes: a mass
on the far-source plane cannot deflect that source's light.  As
$\chi_{s1}\to\chi_l$ the geometric factor tends to one and the two
galaxies act as one plane, with the deflections in the ratio of
$\sigma^2$.

\textbf{(c)} At $(0.5,1.0,2.0)$ the geometric factor is $0.570$.  For
$\sigma_1/\sigma_L=0.4$ the coupling deflection is $9.1\%$ of the main
lens's; for $0.6$ it is $20.5\%$.

\emph{[Attested: \solnlink{14_Multiplane_DSPL/problems\_14.wl}{problems\_14.wl},
checks 6a--6d.]}

% ----- Exercise 14.7 -----
\subsection*{Exercise~\ref{ex:mpi_shape_shear}: What Cancels in $\etamea$}

\textbf{(a)} Write
$\thetaE^{(i)}=[4\pi\sigma_v^2/c^2]\,f(q,\varphi)\,D_{l,si}/D_{si}$.  The
mass factor and $f$ are properties of the deflector, identical for $i=1,2$,
so
$\etamea=(D_{l,s1}/D_{s1})/(D_{l,s2}/D_{s2})=\Brat$ for any $f$.
The argument never used the specific form of $f$, only that the shape of
the deflector is independent of the source redshift.  It also shows that
the ratio is convention-free.

\textbf{(b)} Let the lens-plane mass consist of the SIE plus a perturber
at $z_l$ that produces, near the lens, a constant shear.  Reduced to
plane $i$ the total deflection is
\begin{equation*}
  \vecalpha^{(i)}(\vectheta)=\frac{D_{l,si}}{D_{si}}
  \Big[\hat{\vecalpha}_{\rm SIE}(\vectheta)+\Gamma\vectheta\Big],
  \qquad
  \Gamma=\begin{pmatrix}\gamma_1&\gamma_2\\ \gamma_2&-\gamma_1\end{pmatrix},
\end{equation*}
with one efficiency per plane because both components sit at $z_l$.
Hence $\vecalpha^{(2)}=\vecalpha^{(1)}/\Brat$ exactly, and the reduced
shear itself scales as $\gamma^{(2)}=\gamma^{(1)}/\Brat$.  A model that
gives each plane its own reduced shear, related by the same factor
$1/\Brat$ as the SIE normalisation, reproduces both image systems with
$\etamea=\Brat$, where $\etamea$ is the ratio of the fitted $\thetaE$
parameters.  The radii of the critical curves do not share this exact
scaling: the SIE deflection is of degree zero in $\vectheta$ and the
shear term of degree one, so the critical curve on plane 2 is not a
rescaled copy of the one on plane 1.  If instead a fit forces one shear value on both
planes, it is mis-specified, and the resulting ratio of Einstein radii
is only approximately $\Brat$.

\textbf{(c)} A perturber at $\chi_{s1}<\chi_p<\chi_{s2}$ lies behind
source 1, so its efficiency towards plane 1 is zero, while towards plane
2 it is $(\chi_{s2}-\chi_p)/\chi_{s2}>0$.  It deflects the far source
only.  Its effect is not a common multiple of the main lens's
efficiencies and leaves a genuine plane-to-plane residual in $\etamea$;
the same holds, less dramatically, for any perturber not at $z_l$.

\emph{[Attested: \solnlink{14_Multiplane_DSPL/problems\_14.wl}{problems\_14.wl},
checks 7a--7d; \texttt{multiplane\_dspl.wl} check 4h.]}

% ----- Exercise 14.8 -----
\subsection*{Exercise~\ref{ex:mpi_two_channel}: Two Channels}

From $\ln\Brat=(\gp-1)\ln\etamea$, the total differential, treating
$\etamea$ and $\gp$ as independent, is
\begin{equation*}
  \dd\ln\Brat=(\gp-1)\,\frac{\dd\etamea}{\etamea}+\ln\etamea\,\dd\gp ,
\end{equation*}
which is eq.~\eqref{eq:mpi_twochannel}.  For $0<\etamea<1$,
$\ln\etamea<0$; for $\gp>1$, $\gp-1>0$.  The ratio channel and the slope
channel have opposite signs.  Even at $\gp=2$, where $\Brat=\etamea$, the
slope coefficient $\ln\etamea$ is nonzero: the identity holds at the
isothermal point but is not protected against slope errors.

\emph{[Attested: \solnlink{14_Multiplane_DSPL/problems\_14.wl}{problems\_14.wl},
checks 8a--8d.]}

% ----- Exercise 14.9 -----
\subsection*{Exercise~\ref{ex:mpi_degeneracy}: The Single-Lens Degeneracy}

In parametrisation (b), each likelihood call computes
$\Brat(\Omegam,w)$ and feeds only that number into
eq.~\eqref{eq:mpi_plane2fit}.  So
$\chi^2(\mathbf{p}_{\rm mod},\Omegam,w)=\chi^2(\mathbf{p}_{\rm mod},\Brat(\Omegam,w))$,
and two cosmologies with equal $\Brat$ give identical $\chi^2$ for every
$\mathbf{p}_{\rm mod}$: the likelihood is constant along each level
curve.  On a level curve,
\begin{equation*}
  \dd\Brat=\pderiv{\Brat}{\Omegam}\dd\Omegam+\pderiv{\Brat}{w}\dd w=0
  \;\Longrightarrow\;
  \deriv{w}{\Omegam}\bigg|_{\Brat}
  =-\frac{\partial\Brat/\partial\Omegam}{\partial\Brat/\partial w}
  =-\frac{S_{\Omegam}}{S_w},
\end{equation*}
the factor $\Brat$ cancelling between numerator and denominator.  At the
illustrative geometry $S_{\Omegam}=-0.0357$ and $S_w=-0.0330$, so the
slope is $-1.08$: raising $\Omegam$ (which lowers $\Brat$) must be
compensated by lowering $w$ (which raises it).  Numerically, a step of
$0.01$ in $\Omegam$ along this direction changes $\ln\Brat$ more than
$40$ times less than the same step taken perpendicular to it.

\emph{[Attested: \solnlink{14_Multiplane_DSPL/problems\_14.wl}{problems\_14.wl},
checks 9a--9d.]}

% ----- Exercise 14.10 -----
\subsection*{Exercise~\ref{ex:mpi_error_budget}: Where to Spend Observing Time}

The coefficients of eq.~\eqref{eq:mpi_twochannel} are
$c_\eta=\gp-1=0.60$ and $c_\gamma=\ln\etamea=\ln0.72=-0.3285$.  With
$\sigma_\eta/\etamea=0.01/0.72$ and $\sigma_{\gp}=0.05$:
\begin{equation*}
  \text{ratio channel } c_\eta\frac{\sigma_\eta}{\etamea}=0.00833,\qquad
  \text{slope channel } c_\gamma\sigma_{\gp}=-0.01643 .
\end{equation*}
The variance is
\begin{equation*}
  \Big(\frac{\sigma_\Brat}{\Brat}\Big)^2
  =c_\eta^2\Big(\frac{\sigma_\eta}{\etamea}\Big)^2+c_\gamma^2\sigma_{\gp}^2
  +2\rho\,c_\eta c_\gamma\frac{\sigma_\eta}{\etamea}\sigma_{\gp},
\end{equation*}
and because $c_\gamma<0<c_\eta$ with $\rho>0$, the cross term is
negative, $-2.190\times10^{-4}$.  The total is
$\sigma_\Brat/\Brat=1.10\%$; without the correlation it would be
$1.84\%$, so the positive correlation reduces the error by $40\%$.

The slope channel carries $3.9$ times the variance of the ratio channel.
Better astrometry shrinks only the smaller, ratio channel; pinning the
slope attacks the dominant term.  Stellar kinematics, the radial
thickness of extended arcs and a detected radial or central image are
all ways to do that (\S\ref{sec:mpi_degeneracy}).  For $|S_w|=0.0330$
the result maps to $\sigma_w=0.0110/0.0330=0.33$ from this single lens,
at fixed $\Omegam$.  With $\Omegam$ free, one lens leaves $w$
unconstrained along the degeneracy of
Exercise~\ref{ex:mpi_degeneracy}.

\emph{[Attested: \solnlink{14_Multiplane_DSPL/problems\_14.wl}{problems\_14.wl},
checks 10a--10g and 9a--9d.]}
